\lysh{12631}{2}

\begin{alist}
\item Η κλίση της ευθείας $(\varepsilon_1)$ είναι $\alpha_1 = \dfrac{3}{4}$ και η ευθεία $(\varepsilon_2)$ είναι παράλληλη στην $(\varepsilon_1)$, οπότε έχει κλίση $\alpha_2 = \dfrac{3}{4}$.

\item Η ευθεία $(\varepsilon_2)$ έχει εξίσωση: $y = \dfrac{3}{4}x + \beta$ και διέρχεται από το σημείο $A(4, 1)$, οπότε οι συντεταγμένες του επαληθεύουν την εξίσωση της ευθείας, δηλαδή:
\[1 = \dfrac{3}{4} \cdot 4 + \beta \Leftrightarrow 1 = 3 + \beta \Leftrightarrow \beta = -2.\]

Άρα η ευθεία $(\varepsilon_2)$ έχει εξίσωση: $y = \dfrac{3}{4}x - 2$.

\item Η ευθεία τέμνει τον άξονα $y'y$ στο σημείο $(0, -2)$, αφού για $x = 0$ βρίσκουμε
$y = \dfrac{3}{4} \cdot 0 - 2 = -2$ 
και τον άξονα $x'x$ στο σημείο $\left(\dfrac{8}{3}, 0\right)$, αφού για $y = 0$ έχουμε:
\[0 = \dfrac{3}{4}x - 2 \Leftrightarrow 2 = \dfrac{3}{4}x \Leftrightarrow 8 = 3x \Leftrightarrow x = \dfrac{8}{3}\]
\end{alist}

